\section{Introduction}

Reproducibility is a foundation of cumulative computational research, yet the
term describes several distinct goals. A result may be repeatable in the
original environment, reproducible from shared code and data, or robust to
reasonable changes in software and hardware. Treating these goals as a single
binary property can conceal important failure modes \citep{peng2011,stodden2016}.
For example, a script may execute today but fail after an indirect library is
updated, or it may run successfully while producing a materially different
estimate because a random seed was not recorded.

Practical guidance emphasizes version control, scripted analyses, preserved raw
data, and recorded intermediate results \citep{sandve2013}. The FAIR principles
similarly stress that digital research objects should be findable, accessible,
interoperable, and reusable \citep{wilkinson2016}. These practices improve the
conditions for reuse, but they do not by themselves quantify how resilient a
pipeline is to environmental change. A robustness test is therefore a useful
companion to a reproducibility checklist.

This study presents a small, transparent framework for stress-testing
computational pipelines. We ask three questions:

\begin{enumerate}
  \item How strongly does dependency specification affect execution success?
  \item Do environment capture and random-seed control address the same failure
        mode or complementary ones?
  \item Which compact set of measures communicates pipeline robustness without
        obscuring the underlying failure mechanisms?
\end{enumerate}

The resulting manuscript also demonstrates the expected structure of a general
Taylor \& Francis submission, including author affiliations, an abstract,
keywords, equations, figures, tables, data and code statements, contributor
roles, disclosure, funding, and an author--date reference list.
